\section{Related Work}
\label{sec:related-work}

\subsection{Consensus Strength and Minority Dissent}

Classic research on social influence shows that disagreement with a group is not equivalent across opinion climates. A unanimous majority can alter even judgments about relatively unambiguous stimuli \cite{asch1956independence}, through both normative pressure to avoid social rejection and informational reliance on others \cite{deutsch1955normative}. More directly, the influence of majority and minority sources varies with the level of consensus surrounding their positions \cite{martin2002consensus}, and meta-analytic evidence indicates that smaller deviant factions face stronger rejection \cite{tata1996proportionate}. These findings motivate treating consensus strength as a substantive moderator rather than merely a descriptive property of a thread. Accordingly, we distinguish high-agreement posts, in which YA or NA judgments constitute at least 80\% of verdict-bearing comments, from low-agreement posts, in which the two judgment classes each constitute roughly half. Throughout, \emph{minority} denotes a numerical minority of expressed judgments within a post, rather than a demographic minority or an enduring identity group.

Research on reactions to deviants further suggests that negative evaluation and communication can coexist. In Schachter's foundational study, groups directed communication toward a deviant member in an effort to restore uniformity, yet ultimately rejected the deviant when disagreement persisted \cite{schachter1951deviation}. A later replication recovered the rejection effect but found weaker evidence for concentrated communication \cite{wesselmann2014revisiting}, underscoring the need to test these processes in other settings. Subsequent work has similarly characterized dissent as both a threat to group uniformity and a potential resource for group functioning \cite{jetten2014deviance}. Minority views can promote divergent processing \cite{nemeth1986differential}, exert indirect influence despite limited public acceptance \cite{wood1994minority}, and increase a group's cognitive complexity even while their advocates are socially rejected \cite{curseu2012socially}. This literature therefore anticipates a possible tension between the social cost borne by dissenters and the conversation their dissent generates. It has not, however, established how that tension varies between strong and divided consensus in large, naturally occurring online moral discussions.

\subsection{Evaluative Sanction and Conversational Engagement Online}

On social platforms, votes are not passive measurements of independent preferences. Early ratings can causally shape later evaluations \cite{muchnik2013social,glenski2017rating}, and negative community feedback can influence users' subsequent production and evaluation behavior \cite{cheng2014community}. On Reddit in particular, voting simultaneously expresses evaluation, enforces local norms, and affects content ranking and visibility \cite{graham2021sociomateriality}. A comment's final score should thus be understood as a socially and technologically mediated evaluative outcome, rather than a direct count of agreement alone.

Evaluative and conversational responses also capture different forms of engagement. Likes, shares, and comments respond to different content properties and should not be collapsed into a single engagement construct \cite{tenenboim2022comments}. Likewise, replying, disliking, and flagging are behaviorally distinct reactions to objectionable comments \cite{kalch2017replying}, and disagreement can attract attention without producing endorsement or sharing \cite{dutceac2022cueing}. Most directly, Davis and Graham show that downvoted Reddit comments can receive more direct replies and larger reply trees than comparable upvoted comments, describing the combination as an ``attention reward'' accompanying the emotional cost of negative ratings \cite{davis2021emotional}. Our study builds on this result but changes the inferential target: rather than using a low score as a proxy for norm deviation, we identify whether a comment's substantive moral judgment agrees with the dominant judgment of its post. This permits us to test whether the sanction--engagement split is specific to dissent under strong consensus, whether it remains under divided judgment, and whether minority status primarily affects entry into a reply exchange or its subsequent escalation. We additionally evaluate the aggregate score of the resulting subtree, while treating this measure as reception of the conversation rather than evidence of deliberative quality.

Research on controversy provides a complementary, thread-level perspective. Controversy can increase conversational interest while also creating discomfort, producing a non-monotonic relationship with discussion \cite{chen2013controversy}. On Reddit, mixed positive and negative feedback and the structure of early replies can predict whether a post will become controversial \cite{hessel2019brewing}. These definitions concern controversy around an item, whereas our comparison concerns the distribution of explicit YA and NA judgments and the relative treatment of speakers within the same thread. Prior studies of conversation trees further show that a large thread can arise either from broad one-time participation or from repeated exchanges among a smaller set of users \cite{backstrom2013threads,choi2015reddit}. This distinction motivates our separation of reply entry from escalation: a minority comment may be unusually likely to trigger a confrontation without changing how far an exchange develops once one has begun.

\subsection{Moralized Dissent, Certainty, and Selective Expression}

Moral disagreement may intensify these dynamics because negative and positive moral evaluations are not simple mirror images. Blame tends to be more differentiated and, in many settings, more extreme than praise \cite{guglielmo2019asymmetric}; proscriptive morality also places stricter demands on avoiding wrongdoing than prescriptive morality places on performing good actions \cite{janoffbulman2009proscriptive}. Blame is therefore not only a private assessment but also a communicative act that identifies and regulates perceived norm violations \cite{malle2014blame}. Online, moral-emotional language is associated with greater diffusion \cite{brady2017emotion}, and social feedback can reinforce subsequent expressions of moral outrage \cite{brady2021social}. These accounts motivate our comparison between \emph{blaming dissent} (a minority YA judgment in an NA-dominant post) and \emph{defending dissent} (a minority NA judgment in a YA-dominant post). We do not equate defense with praise; rather, the blame literature motivates the narrower hypothesis that oppositely directed moral judgments need not provoke symmetric responses.

Opinion-climate research also predicts selection in who voices dissent and how dissent is expressed. The spiral-of-silence tradition argues that perceiving oneself to be in the minority can inhibit public expression \cite{noelleneumann1974spiral}. Experiments in online forums likewise find selective posting in response to the surrounding opinion climate \cite{yun2011selective}, while anticipated personal attacks and rejection can suppress minority expression in both online and offline settings \cite{neubaum2018sanctions}. Importantly, such inhibition is heterogeneous: people with greater attitude certainty are more willing to express minority positions in unfavorable climates \cite{matthes2010spiral}. Social consensus can itself affect certainty, although its direction depends on whether belonging or distinctiveness is salient \cite{clarkson2013consensus}, and doubtful language receives fewer recommendations than confident language in large-scale online discussions \cite{evans2023doubt}.

These studies inform, but do not identify, the mechanism behind temporal changes in our data. Because Reddit traces reveal expressed comments rather than unspoken opinions, increasing certainty among later minority comments may reflect selective entry or persistence by highly committed dissenters, within-person attitude strengthening, or both. We therefore describe our annotation as \emph{expressed linguistic certainty}, not latent attitude certainty, and interpret temporal changes as compositional patterns consistent with selective expression rather than direct evidence of individual persuasion.

\subsection{Collective Moral Judgment and Institutionalized Consensus in AITA Communities}

AITA has become an important setting for studying everyday morality because it combines first-person narratives, explicit categorical judgments, voting, and threaded discussion. Nguyen et al. \cite{nguyen2022mapping} mapped more than 100,000 AITA dilemmas into fine-grained topics and related those topics to aggregate YA and NA verdicts. Subsequent research has expanded the empirical map of everyday dilemma types and their aggregate evaluations \cite{yudkin2025large}, examined the situational organization of moral disagreement \cite{shi2025moral}, and analyzed demographic variation in AITA judgments as expressions of social norms \cite{decandia2022socialnorms}. Other studies focus on how narrative features relate to verdicts \cite{giorgi2023author} and on the moral rationales expressed in judgment comments \cite{xi2024morality}. Collectively, this literature has primarily asked what kinds of dilemmas appear, how they are narrated and justified, and which collective verdict they receive. We instead shift the unit of analysis to the unequal reception of individual comments that compose---and contest---the collective verdict.

Three recent studies are especially close to our analysis. Mellody models how agreement and dissent are valued as a consensus develops in AITA, finding that the association between dissent and comment score varies with consensus strength \cite{mellody2025consensus}. Its primary analysis restricts the sample to comments with scores of at least one, following the community's stated rule that downvotes should indicate irrelevance; negatively evaluated comments therefore lie outside its principal estimand. Our analysis retains zero and negative scores because evaluative suppression is itself an outcome of interest, and it jointly examines replies, depth, subtree reception, moral direction, expressed certainty, and time. Goglia and Vega model AITA threads as growing interaction networks and relate thread-level disagreement to their structural evolution \cite{goglia2024structure}; our focus is instead which stance attracts replies within a thread and whether minority status affects conversational entry or escalation. Finally, Goglia, Vega, and Gandelli analyze judgment composition and language before and after AITA's 18-hour verdict disclosure \cite{goglia2026influence}. They report post-disclosure shifts but explicitly caution that self-selection and algorithmic visibility preclude a causal interpretation. Consistent with that caution, we distinguish gradual temporal trends from any discontinuity at 18 hours and interpret a post-lock coefficient as a conditional association rather than a verdict effect.

The comparison between AITA and AITAH provides an institutional boundary condition rather than a natural experiment. Subreddit rules and governance can shape participation and local norms \cite{fiesler2018rules}, and even making community rules more salient can change newcomer participation and compliance \cite{matias2019norms}. AITA and AITAH differ not only in whether a formal verdict is assigned after 18 hours but also in size, moderation, membership, and other unobserved characteristics. We therefore use AITAH to test replication under a community without the formal verdict rule, not to identify the causal effect of flair assignment. Bringing these literatures together, our study asks when numerical dissent becomes a locally deviant voice and traces its reception across four distinct channels---evaluation, conversation, language, and time---while using moral direction and institutional context to characterize where the pattern varies.
